% 编译：xelatex experiment-plan.tex（需要 ctex 宏包，编译两遍以解析引用）
\documentclass[11pt,a4paper]{ctexart}
\usepackage[margin=2.5cm]{geometry}
\usepackage{amsmath,amssymb}
\usepackage{booktabs,tabularx}
\usepackage{enumitem}
\usepackage[hidelinks]{hyperref}
\usepackage{xurl}

\title{近似计算跨层误差传播：FPGA 实验方案}
\date{2026年11月12日}

\begin{document}
\maketitle

\begin{abstract}
本方案利用 ZCU102 与 ZCU104 两块 FPGA 开发板，按“器件—电路—框架—应用”四个抽象层次，分别实测误差的产生与传播，并验证跨层组合模型的预测精度。FPGA 在本研究中的定位是误差行为的真实测量与快速验证平台，不用于论证近似计算的功耗收益。文中板卡参数凭记忆给出，实施前请以 ZCU102 用户手册（UG1182）和 ZCU104 用户手册（UG1267）为准。
\end{abstract}

\section{总体安排}

\begin{table}[h]
\centering
\small
\begin{tabularx}{\textwidth}{@{}lXX@{}}
\toprule
阶段 & 内容 & 平台 \\
\midrule
0 & 环境准备与安全检查 & 两块板 \\
1 & 器件层：超频/降压下的时序错误表征 & ZCU104 \\
2 & 电路层：近似乘法器误差表征；近似误差与时序错误的叠加 & ZCU104 \\
3 & 框架层：逐层误差传播测量 & ZCU102 \\
4 & 应用层：端到端精度测量 & ZCU102 \\
5 & 跨层组合模型验证 & ZCU102 + 主机 \\
6 & 两颗芯片对比 & 两块板 \\
\bottomrule
\end{tabularx}
\caption{实验阶段与平台分工}
\end{table}

分工理由：ZCU102（XCZU9EG，DSP 约 2520、LUT 约 274K）资源更多，用于主加速器；ZCU104（XCZU7EV）有 PYNQ 官方支持，便于用 Python 控制实验，且把有风险的超频、降压实验与主平台隔离。

\section{阶段 0：环境准备与安全检查}

\begin{enumerate}[leftmargin=*]
  \item \textbf{工具链}：安装与 PYNQ 镜像匹配的 Vivado/Vitis 版本；ZCU104 烧写 PYNQ 官方镜像，ZCU102 使用 PetaLinux 或社区 PYNQ 镜像。
  \item \textbf{软件依赖}：主机端安装 PyTorch、AdaPT（近似 DNN 仿真）、NumPy/Pandas；获取 EvoApprox8b 近似乘法器库（Verilog 与 C 模型）。
  \item \textbf{温度监控}：通过片上 SYSMON 读取结温，每次测量都记录温度；确认风扇工作正常。
  \item \textbf{电源记录}：通过 PMBus 读出并保存 VCCINT 等电源轨的默认设置，写成一键恢复脚本。
  \item \textbf{安全规则}：
  \begin{itemize}
    \item 只降压，不升压；电压每步调整后先读回确认。
    \item 结温超过设定上限（如 85\,$^\circ$C）立即停止并恢复默认频率。
    \item 在 Zynq UltraScale+ 上，PL 的 VCCINT 与 PS 电源轨分开（以 UG 为准），因此降低 PL 电压通常不会导致 ARM 端崩溃；但仍需准备断电重启的恢复流程。
  \end{itemize}
  \item \textbf{数据管理}：所有测量结果按统一格式保存为 CSV（见第~\ref{sec:data} 节），原始数据不放进 git 仓库，只提交处理脚本和汇总结果。
\end{enumerate}

\section{阶段 1：器件层——时序错误表征（ZCU104）}

\subsection{测试电路设计}
\begin{enumerate}[leftmargin=*]
  \item \textbf{被测单元}：一个精确的 8 位乘法器，强制用 LUT 实现（综合属性 \texttt{use\_dsp = "no"}），使其关键路径落在可控的逻辑上。
  \item \textbf{时钟}：被测单元的输入、输出寄存器由 MMCM 产生的测试时钟驱动；MMCM 开启动态重配置（DRP），可由 ARM 端在运行时修改频率。
  \item \textbf{激励}：输入向量存放在 BRAM 中，由 ARM 端写入。
  \item \textbf{参考结果}：在安全频率下先运行一遍，或由 ARM 端软件计算，作为正确结果。
  \item \textbf{结果采集}：把每个输出与参考结果的异或（错误位图）写入 BRAM，由 ARM 端读回。
  \item \textbf{控制接口}：AXI-Lite 寄存器，负责启动、停止、设定重复次数。
\end{enumerate}

\subsection{激励设计}
时序错误不仅取决于当前输入，还取决于\textbf{上一个输入到当前输入的跳变}。8 位乘法器的所有输入对共 $2^{16}=65536$ 种，所有“前一输入—当前输入”组合有 $2^{32}$ 种，无法穷举。采用两种激励：
\begin{itemize}[leftmargin=*]
  \item \textbf{双向量测试}：每个输入前插入一个固定输入（如全 0），覆盖全部 65536 种当前输入，便于与静态误差表对照。
  \item \textbf{真实数据序列}：从 DNN 推理中截取真实的权重与激活序列，反映实际工作时的跳变分布。
\end{itemize}

\subsection{频率扫描}
\label{sec:freq}
\begin{enumerate}[leftmargin=*]
  \item 记录 Vivado 时序报告中的标称频率 $f_0$ 和关键路径裕量。
  \item 从 $f_0$ 开始，以 5\,MHz 为步长提高频率，直到错误率超过 50\% 或电路失效。
  \item 每个频率点重复 $R=10$ 次，估计每个输出位的出错概率及其置信区间。
  \item 同时记录温度和电压。
\end{enumerate}

\subsection{电压扫描（频率扫描完成后再做）}
\begin{enumerate}[leftmargin=*]
  \item 固定频率为 $f_0$，通过 PMBus 将 VCCINT 从标称值以 10\,mV 为步长降低。
  \item 每步重复第~\ref{sec:freq} 节的测量，直到出现错误后再降 3--5 步或电路失效为止。
  \item 测量结束后执行恢复脚本。
\end{enumerate}

\subsection{输出}
每个输出位的出错概率表：
\begin{equation}
P(\text{第 } k \text{ 位出错} \mid f,\ V,\ T,\ \text{激励类型}),
\end{equation}
即论文 \cite{vallero2016} 中技术域 SER 库在近似计算中的实测对应。

\section{阶段 2：电路层——近似乘法器误差表征（ZCU104）}

\subsection{静态近似误差}
\begin{enumerate}[leftmargin=*]
  \item 从 EvoApprox8b 选取 4--6 个误差程度不同的 8 位近似乘法器，加上精确乘法器作为对照。
  \item 在安全频率下穷举 65536 种输入，读回结果，得到误差函数 $e(a,b)$。
  \item 与 C 模型逐比特比对，确认硬件实现正确。
  \item 计算标准指标：错误率 ER、MED、MRED、偏置 $\mathbb{E}[e]$、方差 $\operatorname{Var}[e]$。
  \item \textbf{以数据为条件的误差矩}：用真实 DNN 的权重、激活直方图对 $e(a,b)$ 加权，得到 $\mathbb{E}[e \mid \text{数据分布}]$ 和 $\operatorname{Var}[e \mid \text{数据分布}]$，并与均匀分布下的结果对比。
\end{enumerate}

\subsection{近似误差与时序错误的叠加（关键实验）}
\begin{enumerate}[leftmargin=*]
  \item 让每个近似乘法器在阶段 1 的频率/电压扫描点上运行。
  \item 测得总误差 $e_{\text{total}}(a,b)$，分解为近似误差 $e_{\text{approx}}$ 与时序错误 $e_{\text{timing}}$。
  \item 检验叠加关系：
  \begin{equation}
  \mathbb{E}[e_{\text{total}}] \overset{?}{=} \mathbb{E}[e_{\text{approx}}] + \mathbb{E}[e_{\text{timing}}],\qquad
  \operatorname{Var}[e_{\text{total}}] \overset{?}{=} \operatorname{Var}[e_{\text{approx}}] + \operatorname{Var}[e_{\text{timing}}].
  \end{equation}
  \item 注意：近似乘法器的关键路径通常比精确乘法器短，同一频率下时序错误可能更少。需按各自的时序裕量对齐后再比较。
\end{enumerate}

\section{阶段 3：框架层——逐层误差传播测量（ZCU102）}

\subsection{加速器设计}
\begin{enumerate}[leftmargin=*]
  \item \textbf{计算阵列}：INT8 脉动阵列（如 $16\times16$），用 HLS 或 RTL 实现；累加器位宽可配置（如 24 位、32 位）。
  \item \textbf{近似乘法器的实现}，两种方案：
  \begin{itemize}
    \item \textbf{查找表方案（推荐先做）}：把任意近似乘法器存成 $65536 \times 16$ 位的查找表（约 1\,Mb，放在 BRAM 中），运行时替换表内容即可切换乘法器。误差行为与真实电路逐比特一致，但不反映真实电路的时序。
    \item \textbf{真实电路方案}：例化若干种近似乘法器，用选择器按网络层切换。占用面积更大，但能与阶段 1 的超频实验结合。
  \end{itemize}
  \item \textbf{逐层卸载}：ARM 端按网络层依次调用加速器，每层输出经 DMA 读回，因此天然可以拿到每一层的中间结果。
\end{enumerate}

\subsection{测量步骤}
\begin{enumerate}[leftmargin=*]
  \item 对同一批输入，分别用精确乘法器和近似乘法器运行，得到每层输出的误差 $\Delta y^{(l)}$。
  \item 统计每层误差的均值、方差，以及同一层内不同输出之间的协方差。
  \item \textbf{单层注入}：只在第 $l$ 层使用近似乘法器，其余层精确，测量误差在后续各层中的放大或衰减（对应论文中只运行到层边界的 IRS 模式）。
  \item \textbf{全层注入}：所有层都用近似乘法器，测量误差累积。
  \item \textbf{验证点积公式}：对每层比较实测值与解析预测
  \begin{equation}
  \mathbb{E}[\Delta y] = N\mu,\qquad
  \operatorname{Var}[\Delta y] = N\sigma^2 + \sum_{i\neq j}\operatorname{Cov}(\varepsilon_i,\varepsilon_j),
  \end{equation}
  重点检验：忽略协方差项时预测偏差有多大；BatchNorm、ReLU、残差连接对误差的影响。
\end{enumerate}

\section{阶段 4：应用层——端到端精度（ZCU102）}
\begin{enumerate}[leftmargin=*]
  \item \textbf{模型}：先用 CIFAR-10 上的 ResNet-20，再扩展到 MobileNetV2；INT8 量化（训练后量化或量化感知训练）。
  \item \textbf{数据}：CIFAR-10 测试集全集；之后可扩展到 ImageNet 子集（如 1 万张）。
  \item \textbf{指标}：Top-1 精度下降；可选：输出置信度的校准误差。
  \item \textbf{对比}：同一配置下比较三种结果——板上实测、AdaPT 软件仿真、解析模型预测。查找表方案下前两者应逐比特一致，可作为正确性检查。
\end{enumerate}

\section{阶段 5：跨层组合模型验证}
\begin{enumerate}[leftmargin=*]
  \item \textbf{组合模型}：器件层错误概率表（阶段 1）$\rightarrow$ 电路层误差矩（阶段 2）$\rightarrow$ 框架层逐层矩传播（阶段 3 的解析模型）$\rightarrow$ 应用层精度映射（用少量实测点标定）。
  \item \textbf{预测}：给定（频率，电压，各层乘法器配置），预测最终精度。
  \item \textbf{验证}：在至少 20 个未参与标定的配置上实测，报告预测误差（平均绝对误差、最大误差）与耗时对比。这对应论文中用全程故障注入验证贝叶斯模型的做法。
  \item \textbf{消融}：去掉协方差项、改用均匀分布下的误差矩、去掉器件层，分别观察预测误差的变化，量化每个因素的贡献。
  \item \textbf{设计探索}：用组合模型做敏感度分析，选出最敏感的层保持精确，其余层近似；在板上验证选择的效果。
\end{enumerate}

\section{阶段 6：两颗芯片对比}
把阶段 1、2 的测试电路分别部署在两块板上，在相同频率、电压、温度下比较时序错误的起始频率和各位出错概率。只有两个样本，结果仅作为芯片差异的示例，不做统计推断。

\section{数据格式}
\label{sec:data}
\begin{table}[h]
\centering
\small
\begin{tabularx}{\textwidth}{@{}lX@{}}
\toprule
文件 & 字段 \\
\midrule
器件层 & 板卡, 乘法器, 频率, 电压, 温度, 激励类型, 位序号, 出错次数, 总次数 \\
电路层 & 板卡, 乘法器, 频率, 电压, $a$, $b$, 输出, 参考输出 \\
框架层 & 模型, 层号, 乘法器配置, 输入编号, 误差均值, 误差方差, （可选）协方差摘要 \\
应用层 & 模型, 配置, 频率, 电压, Top-1, 来源（板上/仿真/预测） \\
\bottomrule
\end{tabularx}
\caption{测量数据格式}
\end{table}

\section{时间安排（参考）}
\begin{table}[h]
\centering
\small
\begin{tabular}{@{}ll@{}}
\toprule
周次 & 内容 \\
\midrule
第 1 周 & 阶段 0；测试电路与 MMCM 动态重配置调通 \\
第 2--3 周 & 阶段 1 频率扫描；阶段 2 静态误差表征 \\
第 4 周 & 阶段 2 叠加实验；阶段 1 电压扫描 \\
第 5--7 周 & 阶段 3 加速器（查找表方案）与逐层测量 \\
第 8 周 & 阶段 4 端到端精度 \\
第 9--10 周 & 阶段 5 组合模型验证与消融；阶段 6 \\
\bottomrule
\end{tabular}
\caption{参考时间表}
\end{table}

\section{风险与对策}
\begin{itemize}[leftmargin=*]
  \item \textbf{时序错误不可复现}：受温度和电压噪声影响。对策：固定温度区间、多次重复、报告置信区间。
  \item \textbf{超频时 MMCM 或电路整体失效}，而非出现位错误。对策：只让被测单元所在的时钟域超频，控制逻辑保持安全频率。
  \item \textbf{降压导致板卡无法启动}。对策：先完成频率扫描再做降压；准备恢复脚本和断电重启流程。
  \item \textbf{FPGA 布线导致时序错误分布与 ASIC 不同}。对策：明确说明结论适用于 FPGA；用固定布局约束保证多次实验一致。
  \item \textbf{查找表方案不反映真实时序}。对策：只用于框架层和应用层的误差传播验证，时序相关结论来自阶段 1、2 的真实电路。
\end{itemize}

\begin{thebibliography}{9}
\bibitem{vallero2016} A.~Vallero et al. Cross-layer system reliability assessment framework for hardware faults. ITC, 2016. DOI: 10.1109/TEST.2016.7805863.
\bibitem{evoapprox} V.~Mrazek et al. EvoApprox8b: library of approximate adders and multipliers for circuit design and benchmarking of approximation methods. DATE, 2017.
\bibitem{adapt} AdaPT: fast emulation of approximate DNN accelerators in PyTorch. IEEE TCAD, 2022. \url{https://arxiv.org/abs/2203.04071}
\end{thebibliography}

\end{document}
